% Root and D completed independent full literal audits.
\subsection{A small right-zero endpoint behind a singleton long arrow}

\begin{theorem}\label{C:right-zero-small-last-stalk}
Let $C$ be a three-object linear category with all three forward
hom-sets nonempty. Suppose that its last corner is
$R_3=\{1,p,q\}$, with $p^2=p$, $q^2=q$, $pq=q$, and $qp=p$, that
$|C(0,2)|=1$, and that $|C(1,2)|\leq3$.
If $C$ has at most eighteen arrows, then $BC$ is not homotopy
equivalent to $\mathbb CP^2$. There is no projectivity
assumption on either of the other corners. The opposite
statement also holds.
\end{theorem}
\begin{proof}
Let $D$ be the full category on $0,1$ and put $S=C(1,2)$.
The incoming right $D$-act at the last object has a singleton
stalk at $0$ and the left $R_3$-set $S$ at $1$. Write $P,Q$
for the actions of $p,q$ on $S$ and $U=PS=QS$ for their
common image. The two retractions fix $U$ pointwise.

The natural action-category graph has vertex set $U$ and,
for each $s\in S\setminus U$, one edge from $Ps$ to $Qs$.
It can be obtained from the natural labeled graph model by
collapsing all edges labeled $p$, which form disjoint stars.
This is an equivariant weak equivalence: commuting right
maps preserve those stars and send a remaining edge at unit
speed to an edge or collapse it. At object $0$ the graph is
a point. Every map across the unique direction in $D$ takes
the graph at $1$ to that point. Thus these are actual natural
graph diagrams, not merely objectwise homology identifications.
Since $BR_3$ is contractible, the collage has the homotopy model
\[
 X=\operatorname{hocofib}\bigl(\Gamma_{hD}\longrightarrow BD\bigr).
\]

If $|S|\leq3$, the graph at $1$ is either a forest or a disjoint
union of bouquets, with isolated vertices allowed. Indeed
$|U|=1$ gives only loops, $|U|=2$ permits at most one edge,
and $|U|=3$ gives no edges.

In the forest case, replace both graph stalks by their component
sets. The comparison is an objectwise weak equivalence and is
natural for every arrow of $D$. Equivalently, the actual functor
from each $R_3$-action category to its component set has this
property. Collapse the last corner $R_3$ to the trivial monoid,
and replace $S$ by its set of action components. The singleton
long hom-set is unchanged; all composition maps descend.
This defines an actual finite category $C_0$ and a functor
$C\to C_0$ inducing a homotopy equivalence, by the action-bar
homotopy-pushout comparison. Its arrow count is
\[
 |C_0|=|C|-2-\bigl(|S|-|\pi_0\Gamma|\bigr)\leq16.
\]
Theorem~\ref{C:CP2-lower17} excludes this case.

In the remaining case every edge is a loop, so $Ps=Qs$ for
all $s$. The common retraction $P=Q:S\to U$ commutes with
the right action of the middle monoid. It also commutes with
the cross maps to the singleton stalk at $0$. The component
collapse therefore gives an actual functor $q:C\to C_0$,
where the last corner is trivial and its incoming hom-sets
are the singleton and $U$. There is an actual section
$j:C_0\to C$: send the last identity to the identity of
$R_3$, and include the core states $U$ among the arrows.
The singleton long hom-set guarantees that this preserves
all two-cross-arrow compositions. Hence $qj=1$ and
$|C_0|\leq|C|-2\leq16$.

Double each graph loop, fixing its vertex and every point
stalk. This is a strict natural graph map, since all graph
actions send an oriented edge to one edge or to a vertex.
It induces a self-map of $X$, acting as the identity on $BD$.
In the cellular action-bar chain model, graph vertices have
eigenvalue $1$ and graph edges have eigenvalue $2$; the graph
boundary is zero because every edge is a loop. Consequently
the entire cofiber chain complex splits into its component
complex and its loop complex. In every cohomological degree
the only eigenvalues of this self-map are $1,2$, and the
eigenvalue-$1$ subspace is exactly the image of
$q^*:H^*(BC_0;\mathbb Q)\to H^*(BC;\mathbb Q)$.

Suppose now that $BC\simeq\mathbb CP^2$. Its one-dimensional
second rational cohomology is an eigenspace. If its eigenvalue
is $2$, the square of its generator has eigenvalue $4$, which
does not occur in degree four, so that square is zero. This
contradicts the projective-plane ring.

If the eigenvalue is $1$, the generator pulls back from $BC_0$.
Its nonzero square forces both $H^2(BC_0;\mathbb Q)$ and
$H^4(BC_0;\mathbb Q)$ to be nonzero. Since $BC_0$ is an actual
homotopy retract of $BC$, its fundamental group is trivial and
its integral homology is a direct summand of that of
$\mathbb CP^2$. In degrees two and four the nonzero summands
are therefore the whole copy of $\mathbb Z$; there is no
homology in the other positive degrees. The section $Bj$ is
an integral homology equivalence between simply connected
CW spaces and hence a homotopy equivalence. Thus
$BC_0\simeq\mathbb CP^2$, again contradicting its at most
sixteen arrows by Theorem~\ref{C:CP2-lower17}. This proves the theorem.
\end{proof}

\begin{corollary}\label{C:small-last-stalk-seventeen-profile}
There is no seventeen-arrow projective-plane category with
linear profile $(m,n,3;x,y,1)$ satisfying $m\leq2$ and
$y\leq3$. The opposite profiles are also excluded.
\end{corollary}
\begin{proof}
For an ordinary monoid of order three, either its trivial right
rational module is projective or it is the right-zero monoid
$R_3$ above. In the latter case apply the theorem. In the former,
both extreme trivial modules are projective, since $m\leq2$;
Theorem~\ref{C:semisimple-endpoints-seventeen} excludes the
profile. This argument uses the actual projectivity alternatives
and does not infer finite CW type from a rational quotient.
\end{proof}
